\documentclass{article}
\usepackage{geometry}
\geometry{a4paper,scale=0.8}
\usepackage[utf8]{inputenc}

\usepackage{amsmath, amssymb}
\usepackage{amsthm}
\usepackage{enumitem}

\newtheoremstyle{breakstyle}
    {5pt}
    {10pt}
    % {\normalfont}
    {\itshape}
    {0pt}
    {\bfseries}
    {\newline\indent}
    {0pt}
    {\thmname{#1}\thmnumber{ #2}\thmnote{ \textbf{[#3]}}}
\theoremstyle{breakstyle}
\newtheorem{exercise}{Exercise}[section]
\newtheorem{remark}{Remark}[section]

\newenvironment{solution}
  {\begin{proof}[Solution.]\renewcommand{\qedsymbol}{}}
  {\end{proof}}

% \let\ker\relax
% \DeclareMathOperator{\ker}{Ker}
\let\det\relax
\DeclareMathOperator{\det}{Det}

\title{Chapter 1: Finite reflection groups}
\author{Flandre}
\date{\today}

\begin{document}
\maketitle

\section{Reflections}

\begin{exercise}[Trivial]
The reflections form a single conjugacy class in $\mathcal{D}_m$ when $m$ is
odd, but form two classes when $m$ is even.
\end{exercise}

\begin{exercise}[Trivial]
Regarding $S_{n}$ in this way as a subgroup of $\mathrm{O}(n, \mathbf{R})$ ,
prove that the transpositions are the sole reflections belonging to $S_{n}$ .
\end{exercise}

\setcounter{section}{2}
\section{Positive and simple systems}

\begin{exercise}[Trivial]
If $\Phi$ has rank 2, prove that $W$ is a dihedral group. [This will be easier
to do after Theorem 1.5.]
\end{exercise}

\begin{exercise}
Find simple systems for the various groups described in 1.1, taking for $\Phi$ in each case a convenient set of vectors (not necessarily unit vectors).

\begin{solution}
    \begin{enumerate}[label=({\roman*})]
        \item $\left( I_2\left( m \right) ,m\ge 3 \right) $:

            Take
            $$
            \Phi := \left\{ \pm \alpha_j \mid 0 \le j \le m-1,\alpha_j :=\left( \cos j\theta, \sin j\theta \right)  \right\} 
            $$
            and we have simple system
            $$
            \Delta := \{\alpha_0, \alpha_{m-1}\}
            $$
        \item $\left( A_{n-1},n\ge 2 \right) $:

            Take
            $$
            \Phi:=\left\{ \varepsilon_{i}-\varepsilon_{j}\mid i\neq j \right\} 
            $$
            and we have simple system
            $$
            \Delta:=\left\{ \varepsilon_{i}-\varepsilon_{i+1}\mid i\in
            [1,n-1]\cap \mathbb{Z} \right\} 
            $$
        \item $\left( B_{n},n\ge 2 \right) $:

            Take
            $$
            \Phi:=\left\{ \pm \varepsilon_{i},\pm \left(
            \varepsilon_{i}-\varepsilon_{j} \right) \mid 1\le i<j\le n \right\}
            $$
            and we have simple system
            $$
            \Delta:=\left\{
            \varepsilon_1-\varepsilon_2,\varepsilon_2-\varepsilon_3,\ldots
        ,\varepsilon_{n-1}-\varepsilon_{n},\varepsilon_{n} \right\} 
            $$
        \item $\left( D_{n},n\ge 4 \right) $:

            Take
            $$
            \Phi:=\left\{ \pm \left( \varepsilon_{i}\pm \varepsilon_{j} \right)
            \mid 1\le i<j\le n\right\} 
            $$
            and we have simple system
            $$
            \Delta:=\left\{
            \varepsilon_1-\varepsilon_2,\varepsilon_2-\varepsilon_3,\ldots
        \varepsilon_{n-1}-\varepsilon_{n},\varepsilon_{n-1}+\varepsilon_{n} \right\} 
            $$
    \end{enumerate}
    All of them are easy to examine, and we're done.
\end{solution}

\end{exercise}

\setcounter{section}{4}
\section{Generation by simple reflections}

\begin{exercise}[Trivial]
Let $\Phi$ be a root system of rank $n$ consisting of unit vectors. If $\Psi
\subset \Phi$ is a set of $n$ roots whose mutual angles agree with those between
the roots in some simple system, then $\Psi$ must be a simple system.

\begin{proof}
    \emph{(Use \textbf{Corollary 1.5}.)}
\end{proof}

\end{exercise}

\begin{exercise}[Skipped]
Given a simple system $\Delta$ , no proper subset of the simple reflections can
generate $W$ . [Otherwise find $\alpha \in \Delta$ for which $s_{\alpha}$ is not
needed as a generator of $W$ . Consider $w \in W$ for which $w(-\alpha) \in
\Delta$ .]
\end{exercise}

\begin{exercise}[Skipped]
    If $\beta \in \Pi \backslash \Delta$ , prove that $\operatorname{ht}(\beta)
    > 1$ .
\end{exercise}

\section{The length function}

\begin{exercise}
(a) If $w \in W$ , prove that $\det(w) = (-1)^{n(w)}$ . (b) If $w, w' \in W$ ,
prove that $n(ww') \leq n(w) + n(w')$ and $n(ww') \equiv n(w) + n(w') \pmod{2}$
.

\begin{proof}
    \begin{enumerate}[label=({\alph*})]
        \item 
            Because from \textbf{Lemma 1.6}, we know
            $$
            n(w)\equiv l(w)\pmod{2}
            .$$
            So $\det{w}:=\left( -1 \right) ^{l\left( w \right) }=\left( -1
            \right) ^{n\left( w \right) }$.
        \item 
            \emph{(Again use \textbf{Lemma 1.6}, trivial.)}.
    \end{enumerate}
\end{proof}

\end{exercise}

\begin{exercise}[Skipped]
Taking the simple reflections in $S_{n}$ to be the transpositions $(i, i + 1)$ ,
show that the length of a permutation $\pi$ is the number of 'inversions': the
number of pairs $i < j$ for which $\pi(i) > \pi(j)$.
\end{exercise}

\section{Deletion and Exchange Conditions}

\begin{exercise}
In the Exchange Condition, suppose $\ell(w)=r$ . Prove that the index i in the
conclusion is uniquely determined.

\begin{proof}
    Suppose $\exists 1\le i<j\le r,\text{s.t. }$
    $$
    ws=s_1\ldots \hat{s_{i}}\ldots s_{r}=s_1\ldots \hat{s_{j}}\ldots s_{r}
    .$$
    From which we have
    $$
    \hat{s_{i}}s_{i+1}\ldots s_{j-1}s_{j}=s_{i}s_{i+1}\ldots s_{j-1}\hat{s_{j}}
    ,$$
    i.e.
    $$
    s_{i}\ldots s_{j}=s_{i+1}\ldots s_{j-1}
    $$
    and
    $$
    s=s_1\ldots \hat{s_{i}}\ldots \hat{s_{j}}\ldots s_{r}
    $$
    contradicting the minimality of $r$.
\end{proof}

\end{exercise}

\begin{exercise}[Skipped]
Formulate a 'left-handed' version of the Exchange Condition, under the
hypothesis $\ell(sw) < \ell(w)$.

\begin{solution}
Let $w = s_1 \cdots s_r$ (not necessarily reduced), where each $s_i$ is a simple
reflection. If $\ell(sw) < \ell(w)$ for some simple reflection $s$, then there
exists an index $i$ for which $sw = s_1 \cdots \hat{s}_i \cdots s_r$ (and thus
$w = s s_1 \cdots \hat{s}_i \cdots s_r$, with a factor $s$ exchanged for a
factor $s_i$). In particular, $w$ has a reduced expression beginning with $s$ if
and only if $\ell(sw) < \ell(w)$.

\emph{(I don't wanna prove it. Seemed to be trivial.)}
\end{solution}

\end{exercise}

\section{Simple transitivity and the longest element}

\begin{exercise}
What is $w_{\circ}$ in the case of $S_{n}$, relative to the simple system
$\varepsilon_{1} - \varepsilon_{2}, \ldots, \varepsilon_{n-1} - \varepsilon_{n}$
?

\begin{solution}
    In $S_{n},\Pi=\left\{ s_{\varepsilon_{i}-\varepsilon_{j}}\mid i<j \right\}
    $. To send $\Pi$ to $-\Pi$, the required $w_{0}$ is
    $$
    w_0=\binom{1\ 2\ \ldots \ n}{n\ n-1\ \ldots \ 1}
    $$
    And we know
    $$
    n(w_0)=\frac{n(n-1)}{2}
    .$$
\end{solution}

\end{exercise}

\begin{exercise}[Skipped]
In any reduced expression for $w_{0}$, every simple reflection must occur at
least once.
\end{exercise}

\setcounter{section}{9}
\section{Parabolic subgroups and minimal coset  representatives}

\begin{exercise}[Skipped]
Is there a result analogous to (c) describing minimal representatives of the
double cosets $W_{I}wW_{J}(I,J\subset S)$?
\end{exercise}

\begin{exercise}[Skipped]
When $W = S_{n}$ , prove that each parabolic subgroup of $W$ is isomorphic to a
direct product of symmetric groups.
\end{exercise}

\begin{exercise}[Skipped]
Given $s \neq s'$ in $S$ , set $v := ss'ss' \cdots (m$ factors, where $m$ is the
order of $ss'$ in $W$ ), so also $v = s'ss's \cdots (m$ factors), and $v^2 = 1$
. If $w \in W$ satisfies $\ell(ws) < \ell(w)$ and $\ell(ws') < \ell(w)$ , prove
that $\ell(wv) = \ell(w) - m$ . [Consider $W_I$ , $I = \{s, s'\}$ . Since $w =
(wv)v$ with $v \in W_I$ , it suffices to show that $wv \in W^I$ . Look at the
action on the roots corresponding to $s, s'$.]
\end{exercise}

\section{Poincare polynomials}

\begin{exercise}
When $W = S_{3}$ , use the formula in the proposition to compute $W(t)$
inductively, starting with the fact that $W(t) = 1 + t$ for a group of rank 1.
Do the same for $W = \mathcal{D}_m$ in general.

\begin{proof}
    Notice from
    $$
    \sum_{I\subseteq S} \left( -1 \right) ^{I}\frac{W(t)}{W_{I}(t)}=t^{N}
    $$
    we get
    $$
    \sum_{I\subseteq S} \frac{(-1)^{I}}{W_{I}(t)}=\frac{t^{N}}{W(t)}
    .$$
    Breaking up left hand side, we obtian
    $$
    \frac{(-1)^{S}}{W(t)}+\sum_{I\subsetneq S}
    \frac{(-1)^{I}}{W_{I}(t)}=\frac{t^{N}}{W(t)}
    ,$$
    i.e.
    $$
    \frac{t^{N}-(-1)^{S}}{W(t)}=\sum_{I\subsetneq S} \frac{(-1)^{I}}{W_{I}(t)}
    $$
    \emph{(Rest of the proof is pure repitive work, skipped.)}
\end{proof}

\begin{exercise}[Skipped]
The identity just obtained permits an inductive calculation of $|W|$ when $|S|$
is odd. Suppose for example that $|S| = 3$ , and that the dihedral subgroups
$W_{I}$ are of respective orders 4, 6, 10. What is $|W|$?
\end{exercise}

\end{exercise}

\section{Fundamental domains}

\begin{exercise}
Show how the theorem can be used to solve the Word Problem for W: given a
product of simple reflections, decide whether or not it equals 1 in W.

\begin{solution}
    Pick out an element in the corresponding $C$ of the given $W$ and see if
    it's fixed by such product of simple reflections: if so, then the product is
    equal to $1$.
\end{solution}

\end{exercise}

\end{document}
